%-------------------------------------------------------------
\subsection*{R2V control unit}\label{sec:condition_signal}
%-------------------------------------------------------------

Given the extreme viewpoint disparity between the fixed,
elevated roadside sensors and the moving, ground-level vehicle
cameras, we bridge the two platforms through a single
multi-channel conditioning tensor, which we refer to as the
\emph{R2V control unit}. The unit stacks three physically
complementary channels derived from roadside sensor data---a
\emph{photometric channel} that transfers appearance evidence,
a \emph{geometric channel} that enforces metric scale, and a
\emph{layout channel} that carries instance-level
semantics---into a single tensor that serves as the sole
dynamic conditioning signal throughout vehicle-side
generation. All three channels are obtained from the same
cross-platform projection pipeline, so that they are
spatially co-registered to the target vehicle camera at every
frame.

\textbf{Cross-platform projection pipeline.}
Underlying all three channels is a chain of rigid transforms
that links the static roadside frame to the moving
ego-vehicle frame. By combining the static roadside-to-world
extrinsic $T_{wr}$ (Table~\ref{tab:lidar_calib}) with the
time-varying ego-vehicle pose $T_{vw}(t)$
(Eq.~\eqref{eq:Tvw}), any point $\mathbf{p}_r$ in the merged
roadside LiDAR scan is mapped to the ego-vehicle coordinate
system at the matching cooperative timestamp, formulated as
%
\begin{equation}
  \label{eq:chain}
  \mathbf{p}_v = T_{vr}(t)\,\mathbf{p}_r,
  \qquad T_{vr}(t) = T_{vw}(t)\,T_{wr}.
\end{equation}
%
Given $\mathbf{p}_v$, projection onto each of the seven
surround-view cameras $i\in\{1,\ldots,7\}$ then proceeds in
two steps. The point is first expressed in the coordinate
system of the $i$-th vehicle camera via its static extrinsic
$T_{c_v v}^{(i)}$, written as
%
\begin{equation}
  \label{eq:cam}
  \mathbf{p}_{c_v}^{(i)} = T_{c_v v}^{(i)}\,\mathbf{p}_v.
\end{equation}
%
Writing $\mathbf{p}_{c_v}^{(i)} = (x_c, y_c, z_c)^\top$, the
normalised pinhole coordinates are
$(\tilde{x},\tilde{y})=(x_c/z_c,\,y_c/z_c)$. The plumb-bob
model then corrects for the non-negligible radial and
tangential lens distortion of the vehicle cameras as
%
\begin{equation}
  \label{eq:distort}
  \begin{bmatrix}x_d\\y_d\end{bmatrix}
  = \bigl(1+k_1r^2+k_2r^4+k_3r^6\bigr)
    \begin{bmatrix}\tilde{x}\\\tilde{y}\end{bmatrix}
  + \begin{bmatrix}
      2p_1\tilde{x}\tilde{y}+p_2(r^2+2\tilde{x}^2)\\
      p_1(r^2+2\tilde{y}^2)+2p_2\tilde{x}\tilde{y}
    \end{bmatrix},
\end{equation}
%
where $r^2=\tilde{x}^2+\tilde{y}^2$ and
$(k_1, k_2, k_3, p_1, p_2)$ are the per-camera distortion
coefficients. The distortion-corrected coordinates are
finally mapped to pixel coordinates by the intrinsic matrix,
%
\begin{equation}
  \label{eq:pixel}
  \begin{bmatrix}u\\v\end{bmatrix}^{\!(i)}
  = \begin{bmatrix}f_x^{(i)}&0&c_x^{(i)}\\0&f_y^{(i)}&c_y^{(i)}\end{bmatrix}
    \begin{bmatrix}x_d\\y_d\\1\end{bmatrix},
\end{equation}
%
where $(f_x^{(i)},\,f_y^{(i)})$ and $(c_x^{(i)},\,c_y^{(i)})$
denote the focal lengths and principal-point offsets of the
$i$-th vehicle camera as defined in Eq.~\eqref{eq:intrinsic}.
For brevity we encapsulate this end-to-end spatial-to-pixel
mapping as $\pi^{(i)}(\mathbf{p}_r)$, and define analogous
mappings $\pi_r^{(j)}(\cdot)$ for each of the four roadside
pinhole cameras $j\in\{1,\ldots,4\}$ by substituting the
corresponding roadside extrinsics and intrinsics.

\textbf{Photometric channel.}
Synthesising a photorealistic driving scene requires more
than a discrete layout of objects; the generator must also
recover global context such as ambient illumination,
road-surface textures, and background architecture. We
capture this low-frequency appearance prior by colouring each
roadside point $\mathbf{p}_r$ with the mean RGB value sampled
from the roadside pinhole cameras that observe it,
%
\begin{equation}
  \label{eq:blur}
  c(\mathbf{p}_r) = \frac{1}{|\mathcal{J}(\mathbf{p}_r)|}
  \sum_{j \in \mathcal{J}(\mathbf{p}_r)}
  \mathbf{I}_r^{(j)}\!\left(\pi_r^{(j)}(\mathbf{p}_r)\right),
  \qquad
  \mathcal{J}(\mathbf{p}_r) = \{j : \pi_r^{(j)}(\mathbf{p}_r) \in \Omega_j\},
\end{equation}
%
where $\mathbf{I}_r^{(j)}$ denotes the $j$-th roadside image,
$\Omega_j$ its valid pixel domain, and
$\mathcal{J}(\mathbf{p}_r)$ the set of roadside cameras whose
projection of $\mathbf{p}_r$ falls inside that domain. The
coloured points are then projected onto each vehicle camera
via $\pi^{(i)}$, yielding a sparse scatter image whose
non-empty pixels encode the illumination and structural
layout observed from the infrastructure. A pixel is left
black when no roadside point projects onto it, or when the
projecting point itself was not observed by any roadside
pinhole camera ($\mathcal{J}(\mathbf{p}_r)=\varnothing$).
Such sparsity is retained deliberately: it conveys the true
field-of-view limits of the roadside array to the generator,
rather than imputing a potentially misleading prior.

\textbf{Geometric channel.}
Without an explicit metric signal, the generator is free to
hallucinate physically implausible scale relationships or to
violate three-dimensional consistency across the seven
vehicle cameras. We supply such a signal as a sparse depth
image, in which, for every pixel, a depth buffer retains the
smallest camera-axis depth among all roadside points that
project to that pixel, given by
%
\begin{equation}
  \label{eq:depth}
  D^{(i)}(u,v) =
  \min_{\mathbf{p}_r:\,\pi^{(i)}(\mathbf{p}_r)=(u,v)}
  \bigl[\mathbf{p}_{c_v}^{(i)}\bigr]_{z_c},
\end{equation}
%
where the minimum is taken over all roadside points
$\mathbf{p}_r$ whose projection under $\pi^{(i)}$ coincides
with pixel $(u,v)$, and $[\cdot]_{z_c}$ denotes the
component along the $z$-axis of the $i$-th vehicle camera.
Raw depth values are retained to preserve the physical
metric scale; a per-frame min--max normalisation
(near$=$white, far$=$black) is additionally recorded for
visualisation. By supplying the three-dimensional scene
topology from the target viewpoint, the geometric channel
grounds 2D pixel synthesis in the underlying scene geometry.

\textbf{Layout channel.}
The elevated, multi-pole roadside infrastructure provides
near-complete coverage of the intersection, enabling the
annotation of every dynamic traffic participant in a single
merged scan---a level of observability unattainable from the
ego-vehicle alone. Projecting these complete annotations into
the vehicle-camera frame reconstructs the full traffic flow
at the intersection, which is a central objective of the R2V
task. Each roadside 3D annotation is parameterised by centre
$\mathbf{c}^R$, half-extents $(l,w,h)$ and yaw $\psi$; its
eight corners are given by
%
\begin{equation}
  \label{eq:corners}
  \mathbf{v}_k^R = \mathbf{c}^R
  + R_\psi\,[\pm l,\,\pm w,\,\pm h]^\top_k,
  \qquad k=1,\ldots,8,
\end{equation}
%
where $R_\psi\!\in\!SO(3)$ is the yaw rotation matrix and
$[\pm l,\pm w,\pm h]^\top_k$ enumerates the eight sign
combinations of the box half-extents. These corners are
mapped to the ego-vehicle frame via Eq.~\eqref{eq:chain} and
then to the image plane via $\pi^{(i)}$. Rendering proceeds
by drawing each box as a solid filled polygon with a
category-specific colour, with correct occlusion handled by
backface culling and a painter's-algorithm depth sort. The
resulting image constitutes an explicit semantic layout of
all annotated traffic participants within the ego-vehicle's
field of view.

\textbf{Temporal segmentation and training units.}
With the three channels defined at every cooperative frame,
the 86 scenes are segmented into training units. Each
cooperative sequence is divided into non-overlapping 29-frame
clips at 10\,Hz that centre on the vehicle's traversal of the
intersection. The clip length is determined by the temporal
compression ratio of the Wan~2.1 video
tokenizer~\cite{wan2pt1}: with a latent temporal resolution
$T_{\mathrm{lat}}\!=\!8$ and a $4\!\times$ temporal
compression, a pixel clip of length
$T_{\mathrm{pix}}=4(T_{\mathrm{lat}}-1)+1=29$ is the shortest
window that maps onto a full-resolution latent tensor. Raw
vehicle-camera frames are downsampled to $1280\!\times\!720$
to match the backbone resolution. For each of the seven
vehicle cameras, every clip yields a synchronised training
triplet consisting of one ground-truth RGB video, one R2V
control unit video formed by stacking the photometric,
geometric and layout channels, and one structured text prompt
describing the scene context and driving direction.
Collectively, these triplets constitute the post-training
data used in the generation model (see \emph{R2V generation
model}).
